% !TEX root = ../../../main.tex
% Block card: catalogue.gramian_gadf -- "Gramian GADF" (Part I, the blocks chapter).
\cardfield{What it does}
Turns a short stretch of one channel into a square picture in which every pixel shows how
two moments of the stretch \emph{differ}: the sister of \blockref{catalogue.gramian_gasf},
whose card explains the idea.

\cardfield{When to use it}
Beside the summation field, when you want to see change rather than level: where the
signal rises or falls relative to other moments. On its own it is mostly for looking; the
trained networks take it through \blockref{catalogue.window_images}.

\cardfield{What goes in, what comes out}
In goes a span of one channel, at most 5{,}000 samples. Out comes a square image with one
row and one column per sample. The image is mirror-reversed across its diagonal (the top
right is the negative of the bottom left), and the diagonal itself is always zero.

\cardfield{Settings that matter}
None. The span you choose is everything.

\cardfield{Watch out for}
The same as the summation field: the span's own lowest and highest values set the scale,
so one spike flattens the rest; a flat span or one missing sample gives an image with no
values; a span over 5{,}000 samples is refused.
